% META: {"id": "TSPE-DERCONV-EX-042", "chapitre": "TSPE-DERIVATION-CONVEXITE", "type_objet": "exercice", "capacites": ["TSPE-DERCONV-C5"], "capacites_codes": ["C5"], "methodes": ["M4"], "parcours": 3, "competences": ["modeliser","raisonner"], "duree_min": 25, "mode_creation": "ex_nihilo", "sources_inspiration": [], "fichier_tex": "chapitres/TSPE-DERIVATION-CONVEXITE/exercices/TSPE-DERCONV-EX-042.tex", "corrige_tex": "chapitres/TSPE-DERIVATION-CONVEXITE/corriges/TSPE-DERCONV-CO-042.tex", "status": "generated"}
% BEGIN-VERIFY
% from sympy import *
% x = symbols('x')
% # f(x) = x^4 - 4x^2 + 1, trouver nb de points d'inflexion et leur nature
% f = x**4 - 4*x**2 + 1
% fpp = diff(f, x, 2)
% assert expand(fpp) == 12*x**2 - 8
% assert solve(fpp, x) == [-sqrt(6)/3, sqrt(6)/3]
% assert fpp.subs(x, 0) == -8  # concave
% assert fpp.subs(x, 2) == 40  # convexe
% f.subs(x, sqrt(6)/3)
% assert simplify(f.subs(x, sqrt(6)/3) - Rational(-11, 9)) == 0
% END-VERIFY
\begin{exercice}{TSPE-DERCONV-EX-042}{3}{25}
On considere $f(x) = x^4 - 4x^2 + 1$ sur $\mathbb{R}$.
\begin{enumerate}
  \item Combien de points d'inflexion la courbe possede-t-elle ? Les determiner.
  \item Esquisser la courbe en indiquant les variations et la convexite.
  \item La courbe coupe-t-elle l'axe des abscisses ? Combien de fois ?
\end{enumerate}

\end{exercice}
